\scanpage{221}
\unsure{Relations|Relation} entre les $4$ complexes, et fonctorialité.

Les \deleted{complexes} algèbres différentielles sont fonctorielles en $(A,M)$. Pour $A$ fixe, $M$ variable, elles transforment somme en produit tensoriel (d'où, pour chacune, application diagonale compatible avec \deleted{$d$ ou $\delta$} la différentielle), et commutent aux \uline{lim} filtrantes.

L'hom can. d'algèbres (fonctoriel en $A,M$)
\diagram{\begin{tikzcd}
\operatorname{Sym}^*M\otimes\Lambda^*M \arrow[r]\arrow[d,equal] & \Gamma^*M\otimes\Lambda^*M\arrow[d,equal]\\
C^*_{DR}\operatorname{Sym}^*M\arrow[r] & C^*_{DR\,\mathrm{p.d.}}(\Gamma^*M)
\end{tikzcd}}
est compatible avec les différentielles $d$. (C'est un isom \deleted{pour} si $A$ de car $0$ (réciproque: vraie si \deleted{$A$} $M$ libre de rang $>0$), donc le calcul de coh. dans $1$ est contenu dans celui de $3$. [Ce dernier se fait en réduisant au cas où $M$ est libre de \unsure{t.f.|d.f.} par Lazard, puis par Künneth au cas $M$ libre de rang $1$, auquel cas c'est trivial].

L'hom can. d'algèbres
\[
\operatorname{Sym}^*M\otimes\Lambda^*M\unsure{-|\longrightarrow}\Gamma^*M\otimes\Lambda^*M
\]
est aussi compatible avec les opérateurs diff. $\delta$. \deleted{\unsure{Comme}} \deleted{\unsure{Comme}} \deleted{\unsure{Alors}} Procédant comme dessus, le calcul de l'homologie de $(4)$ est ramené à celui de $(2)$, qui peut se traiter par la même méthode que pour $(1)$.

La \uline{définition} de l'op $\delta$ dans $(2)$ et $(4)$ peut d'ailleurs se ramener à un cas plus général de Koszul: si $S$ est un anneau comm., $M$ un $S$-module, \deleted{\unsure{not}} et $\varphi:M-S$ une forme
\margin{220}
